%----------
%	FUENTE DE DATOS — corpus IMIO / IRIS
%----------

El corpus empleado en este trabajo ha sido construido en el marco del proyecto
IRIS y está formado por \textbf{7.115 piezas periodísticas} publicadas en prensa
generalista española. A diferencia de los corpus de redes sociales, habituales
en la detección automática de discurso de odio, se trata de textos extensos
—con una longitud mediana de \textbf{5.881 caracteres}— redactados por
profesionales y sometidos a procesos editoriales, lo que sitúa el fenómeno
analizado en el terreno del sesgo sutil antes que en el de la agresión
explícita.

\subsubsection{Composición temporal y por cabecera}

El corpus abarca dos periodos diferenciados, según se recoge en la
Tabla~\ref{tab:iris-corpus-anios}: una serie histórica que cubre de 2017 a 2021
y una campaña de recogida correspondiente a 2024.

\begin{table}[H]
    \centering
    \begingroup
    \small
    \ttabbox[\FBwidth]{
        \caption{Distribución temporal del corpus IRIS}
        \label{tab:iris-corpus-anios}
    }{
        \begin{tabularx}{\textwidth}{
            >{\centering\arraybackslash}X
            >{\centering\arraybackslash}X
            >{\centering\arraybackslash}X
            >{\centering\arraybackslash}X
            >{\centering\arraybackslash}X
            >{\centering\arraybackslash}X
            >{\columncolor{gray!15}\centering\arraybackslash}X
        }
            \hline
            \textbf{2017} & \textbf{2018} & \textbf{2019} & \textbf{2020} & \textbf{2021} & \textbf{2024} & \textbf{Total} \\
            \hline
            976 & 1.199 & 866 & 527 & 565 & 2.982 & 7.115 \\
            \hline
        \end{tabularx}
    }
    \endgroup
\end{table}

Las piezas proceden de siete cabeceras, con una distribución desigual: la
cabecera más representada aporta 3.117 piezas (43,8\,\% del total) y la menos
representada 63 (0,9\,\%). Esta asimetría, propia de un corpus construido con
criterios de relevancia editorial antes que de balanceo estadístico, no afecta
al diseño experimental de este trabajo —que evalúa el acuerdo pieza a pieza—
pero debe considerarse si se pretendiese extraer conclusiones agregadas sobre
el comportamiento comparado de los medios.

\subsubsection{Obtención del texto completo}

La base de datos original recoge los metadatos de cada pieza (fecha, cabecera,
titular, autoría, sección y sus correspondientes codificaciones), pero no su
cuerpo. Dado que el análisis del sexismo lingüístico requiere necesariamente el
texto íntegro, se desarrolló un proceso de extracción automática que recupera el
contenido a partir de la URL de cada pieza. El resultado se incorpora al corpus
en el campo \texttt{contenido\_articulo}, que constituye la entrada de todos los
experimentos descritos en este trabajo.

\subsubsection{Subconjunto de evaluación}
\label{subsec:iris_subconjunto}

De las 7.115 piezas del corpus, \textbf{1.315 cuentan con anotación experta
simultánea en las cinco variables} objeto de estudio, y son las que constituyen
el conjunto de evaluación. De ellas, 1.313 disponen de texto recuperado con
extensión suficiente para el análisis, cifra sobre la que se calculan todas las
métricas del Capítulo~\ref{results_and_discussion}.

Este subconjunto presenta dos particularidades que condicionan el alcance de las
conclusiones:

\begin{itemize}
    \item \textbf{Restricción temporal.} La totalidad de las piezas anotadas
    corresponde a \textbf{2024}. Los resultados obtenidos no pueden extrapolarse
    sin cautela al periodo 2017--2021, cuyas piezas permanecen sin anotar.

    \item \textbf{Cobertura parcial.} El subconjunto anotado representa el
    18,5\,\% del corpus y procede de cinco de las siete cabeceras, con
    predominio de una de ellas (805 piezas, 61,2\,\% del subconjunto).
\end{itemize}

Las 5.800 piezas restantes carecen de anotación de referencia y, por tanto, no
intervienen en el cálculo de las métricas de acuerdo. Su procesamiento resulta
de interés únicamente para el análisis descriptivo y para la prueba de concepto
editorial, escenarios en los que no se dispone de verdad de referencia contra la
que contrastar.
